\documentclass[11pt]{article}
\usepackage[a4paper,margin=18mm]{geometry}
\usepackage[no-math]{fontspec}
\setmainfont{Latin Modern Roman}
\usepackage{amsmath,amssymb,amsthm,xfrac,xspace,hyperref,graphicx,comment}
\excludecomment{DefTralics}

\providecommand{\bibliofont}{\small}
\newcommand{\ie}{i.e.,\xspace}
\newcommand{\eg}{e.g.,\xspace}
\newcommand{\etc}{etc.\xspace}
\newcommand{\cf}{cf.\xspace}
\newcommand{\resp}{resp.\xspace}
\newcommand{\smaller}{\small}
\newcommand{\Subsection}{\subsection}
\newcommand{\Subsubsection}{\subsubsection}
\newcommand{\skpt}{\smallskip}
\emergencystretch=3em
\numberwithin{equation}{section}\input{author-preamble.tex}
\newfontfamily\jepindicatorfont[ItalicFont={Latin Modern Math}]{Latin Modern Math}
\renewcommand{\mathbbmm}[1]{\text{\jepindicatorfont\upshape\char"1D7D9}}
\begin{document}
\begin{center}{\Large Stable item 2036}\end{center}
\paragraph{Source and attribution.}
Serena Guarino Lo Bianco, Domenico Angelo La Manna, Bozhidar Velichkov, \emph{A two-phase problem with Robin conditions on the free boundary}, Journal de l’École polytechnique — Mathématiques 8 (2021), publisher version, DOI 10.5802/jep.139. \href{https://jep.centre-mersenne.org/articles/10.5802/jep.139/}{Publisher article}. Authors' text: \href{https://creativecommons.org/licenses/by/4.0/}{CC BY 4.0}.
\paragraph{Editorial problem boundary.}
Prove only Theorem 1.1(i)--(ii) below: existence of a minimiser for the original two-phase Robin energy, and interior Hölder regularity with a strictly positive lower bound, under all printed \(\beta>0\), smooth bounded domain, exterior finite-perimeter data and positive Sobolev-data assumptions. The complete approximating/limit construction and nondegeneracy core are retained. Free-boundary singular-set and smooth-boundary claims are context only.
\paragraph{Difficulty (editorial): H2.}
Ordinary finite-perimeter compactness does not give a perimeter bound when the Robin boundary weight can approach zero. The author constructs positive-floor approximating minimisers, proves uniform interior Hölder bounds, and passes a carefully chosen optimality condition to the limit. A nonlinear coarea/isoperimetric iteration then forces a positive lower bound, which supplies the missing perimeter compactness and completes the minimiser construction.
\paragraph{Editorial transcription note.}
Original author language, mathematics and ordering are preserved. Publisher front matter is replaced by this independent article wrapper. Bibliography is recovered from matching cached publisher HTML, including its TeX formulas. No new proof or mathematical repair is supplied. Surrounding original results are retained for conservative context and are not extra selected corpus claims.
\clearpage
\input{author-body.tex}
\begin{thebibliography}{99}
\bibitem{bg} D. Bucur \& A. Giacomini - “Shape optimization problems with Robin conditions on the free boundary” , Ann. Inst. H. Poincaré Anal. Non Linéaire 33 (2016) no. 6, p. 1539-1568
\bibitem{bl} D. Bucur \& S. Luckhaus - “Monotonicity formula and regularity for general free discontinuity problems” , Arch. Rational Mech. Anal. 211 (2014) no. 2, p. 489-511
\bibitem{ck} L. A. Caffarelli \& D. Kriventsov - “A free boundary problem related to thermal insulation” , Comm. Partial Differential Equations 41 (2016) no. 7, p. 1149-1182
\bibitem{css} L. A. Caffarelli, M. Soria-Carro \& P. R. Stinga - “Regularity for $C^{1,\alpha }$ interface transmission problems” , 2020
\bibitem{D} H. Dong - “A simple proof of regularity for $C^{1,\alpha }$ interface transmission problems” , 2020
\bibitem{eg} L. C. Evans \& R. F. Gariepy - Measure theory and fine properties of functions , Textbooks in Math. , CRC Press, Boca Raton, FL, 2015
\bibitem{maggi} F. Maggi - Sets of finite perimeter and geometric variational problems. An introduction to geometric measure theory , Cambridge Studies in Advanced Math. , vol. 135 , Cambridge University Press, Cambridge, 2012
\bibitem{tamanini} I. Tamanini - Regularity results for almost minimal oriented hypersurfaces in $\mathbb{R}^n$ , Dipartimento di Matematica dell’Università di Lecce, Lecce, 1984
\bibitem{velectures} B. Velichkov - “Regularity of the one-phase free boundaries” (2019), Lecture notes, available at http://cvgmt.sns.it/paper/4367/
\end{thebibliography}
\end{document}
